% GENERATED FILE. Edit canonical lessons, then run npm run generate:books.
% canonical-source-hash: fnv1a64:7bbf877f822da84c
% canonical-lessons: UR-C427-saiz, UR-C427-dhila, UR-C427-poshak, UR-C427-rumal, UR-C427-sajna

\chapter{Sizes and Outfits}
\label{ch:ur-sizes-and-outfits}
\hlchaptermodality{427}

\begin{chapteropening}
\textbf{By the end of this chapter:} \emph{I can say whether clothes fit and look good.}

Complete the last lesson of chapter 427: I can say whether clothes fit and look good.
\end{chapteropening}

\section[\texorpdfstring{\ur{سائز} (sāiz)}{سائز (sāiz)}]{\ur{سائز} (sāiz) — a size}

\label{lesson:UR-C427-saiz}



\begin{quote}
Before the new one: say the Urdu for to rub, then the Urdu for to polish.
\end{quote}

\subsection*{You'll want to know: \ur{سائز}}
\textbf{\ur{سائز}} — \emph{sāiz} — \textquotedblleft{}a size\textquotedblright{}.

\begin{quote}
\ur{اس} \ur{قمیض} \ur{کا} \ur{سائز} \ur{چھوٹا} \ur{ہے۔}
\end{quote}

\begin{quote}
\emph{is qamīz kā sāiz chhoṭā hai.}
\end{quote}

\begin{quote}
This shirt is a small size.
\end{quote}

Masculine.

\begin{cousinweb}
\textbf{\ur{سائز}} (\emph{sāiz}) is borrowed from English \emph{size}.
\end{cousinweb}

\begin{grammarlens}[title={What you've built}]
One more word for getting around.
\end{grammarlens}

\subsection*{Your turn}
\begin{itemize}
\raggedright
  \item \emph{Say it:} \emph{sāiz}
  \item \emph{Say it:} \emph{sāiz}, once more
  \item \emph{Say it:} \emph{is qamīz kā sāiz chhoṭā hai.} — This shirt is a small size.
  \item \emph{Recall:} say the Urdu for to rub, then the Urdu for to polish, then say \emph{sāiz} again
\end{itemize}

\begin{tcolorbox}[breakable,colback=teal!4,colframe=teal!35!black,title={Before you move on}]
What does \ur{سائز} mean? (\textquotedblleft{}A size\textquotedblright{}.) Say it once more.
\end{tcolorbox}
\section[\texorpdfstring{\ur{ڈھیلا} (ḍhīlā)}{ڈھیلا (ḍhīlā)}]{\ur{ڈھیلا} (ḍhīlā) — loose; slack}

\label{lesson:UR-C427-dhila}



\begin{quote}
Before the new one: say the Urdu for to polish, then the Urdu for a size.
\end{quote}

\subsection*{You'll want to know: \ur{ڈھیلا}}
\textbf{\ur{ڈھیلا}} — \emph{ḍhīlā} — \textquotedblleft{}loose; slack\textquotedblright{}.

\begin{quote}
\ur{یہ} \ur{پتلون} \ur{ڈھیلی} \ur{ہے۔}
\end{quote}

\begin{quote}
\emph{yih patlūn ḍhīlī hai.}
\end{quote}

\begin{quote}
These trousers are loose.
\end{quote}

\begin{cousinweb}
\textbf{\ur{ڈھیلا}} (\emph{ḍhīlā}) is inherited: CDIAL lists the Hindi cousin \emph{ḍhīlā}, \textquotedblleft{}slack, loose\textquotedblright{}, under \emph{ḍhila}, a form it reconstructs (CDIAL 5590).
\end{cousinweb}

\begin{grammarlens}[title={What you've built}]
One more word for getting around.
\end{grammarlens}

\subsection*{Your turn}
\begin{itemize}
\raggedright
  \item \emph{Say it:} \emph{ḍhīlā}
  \item \emph{Say it:} \emph{ḍhīlā}, once more
  \item \emph{Say it:} \emph{yih patlūn ḍhīlī hai.} — These trousers are loose.
  \item \emph{Recall:} say the Urdu for to polish, then the Urdu for a size, then say \emph{ḍhīlā} again
\end{itemize}

\begin{tcolorbox}[breakable,colback=teal!4,colframe=teal!35!black,title={Before you move on}]
What does \ur{ڈھیلا} mean? (\textquotedblleft{}Loose; slack\textquotedblright{}.) Say it once more.
\end{tcolorbox}
\section[\texorpdfstring{\ur{پوشاک} (poshāk)}{پوشاک (poshāk)}]{\ur{پوشاک} (poshāk) — an outfit; a dress}

\label{lesson:UR-C427-poshak}



\begin{quote}
Before the new one: say the Urdu for a size, then the Urdu for loose.
\end{quote}

\subsection*{You'll want to know: \ur{پوشاک}}
\textbf{\ur{پوشاک}} — \emph{poshāk} — \textquotedblleft{}an outfit; a dress\textquotedblright{}.

\begin{quote}
\ur{شادی} \ur{کے} \ur{لیے} \ur{سارا} \ur{نے} \ur{نئی} \ur{پوشاک} \ur{خریدی۔}
\end{quote}

\begin{quote}
\emph{shādī ke liye sārā ne naī poshāk kharīdī.}
\end{quote}

\begin{quote}
Sara bought a new outfit for the wedding.
\end{quote}

Feminine.

\begin{cousinweb}
\textbf{\ur{پوشاک}} (\emph{poshāk}) is usually traced to Persian \emph{poshāk}, \textquotedblleft{}clothing, a dress\textquotedblright{}.
\end{cousinweb}

\begin{grammarlens}[title={What you've built}]
One more word for getting around.
\end{grammarlens}

\subsection*{Your turn}
\begin{itemize}
\raggedright
  \item \emph{Say it:} \emph{poshāk}
  \item \emph{Say it:} \emph{poshāk}, once more
  \item \emph{Say it:} \emph{shādī ke liye sārā ne naī poshāk kharīdī.} — Sara bought a new outfit for the wedding.
  \item \emph{Recall:} say the Urdu for a size, then the Urdu for loose, then say \emph{poshāk} again
\end{itemize}

\begin{tcolorbox}[breakable,colback=teal!4,colframe=teal!35!black,title={Before you move on}]
What does \ur{پوشاک} mean? (\textquotedblleft{}An outfit; a dress\textquotedblright{}.) Say it once more.
\end{tcolorbox}
\section[\texorpdfstring{\ur{رومال} (rūmāl)}{رومال (rūmāl)}]{\ur{رومال} (rūmāl) — a handkerchief}

\label{lesson:UR-C427-rumal}



\begin{quote}
Before the new one: say the Urdu for loose, then the Urdu for an outfit.
\end{quote}

\subsection*{You'll want to know: \ur{رومال}}
\textbf{\ur{رومال}} — \emph{rūmāl} — \textquotedblleft{}a handkerchief\textquotedblright{}.

\begin{quote}
\ur{میری} \ur{جیب} \ur{میں} \ur{ایک} \ur{رومال} \ur{ہے۔}
\end{quote}

\begin{quote}
\emph{merī jeb meṅ ek rūmāl hai.}
\end{quote}

\begin{quote}
There's a handkerchief in my pocket.
\end{quote}

Masculine.

\begin{cousinweb}
\textbf{\ur{رومال}} (\emph{rūmāl}) is built from Persian \emph{rū}, \textquotedblleft{}the face\textquotedblright{}, and \emph{māl}, \textquotedblleft{}wiping\textquotedblright{}: it is usually traced to Persian \emph{rūmāl}.
\end{cousinweb}

\begin{grammarlens}[title={What you've built}]
One more word for getting around.
\end{grammarlens}

\subsection*{Your turn}
\begin{itemize}
\raggedright
  \item \emph{Say it:} \emph{rūmāl}
  \item \emph{Say it:} \emph{rūmāl}, once more
  \item \emph{Say it:} \emph{merī jeb meṅ ek rūmāl hai.} — There's a handkerchief in my pocket.
  \item \emph{Recall:} say the Urdu for loose, then the Urdu for an outfit, then say \emph{rūmāl} again
\end{itemize}

\begin{tcolorbox}[breakable,colback=teal!4,colframe=teal!35!black,title={Before you move on}]
What does \ur{رومال} mean? (\textquotedblleft{}A handkerchief\textquotedblright{}.) Say it once more.
\end{tcolorbox}
\section[\texorpdfstring{\ur{سجنا} (sajnā)}{سجنا (sajnā)}]{\ur{سجنا} (sajnā) — to look good (on someone); to be decorated}

\label{lesson:UR-C427-sajna}



\begin{quote}
Before the new one: say the Urdu for an outfit, then the Urdu for a handkerchief.
\end{quote}

\subsection*{You'll want to know: \ur{سجنا}}
\textbf{\ur{سجنا}} — \emph{sajnā} — \textquotedblleft{}to look good (on someone); to be decorated\textquotedblright{}.

\begin{quote}
\ur{یہ} \ur{پوشاک} \ur{سارا} \ur{پر} \ur{سجتی} \ur{ہے۔}
\end{quote}

\begin{quote}
\emph{yih poshāk sārā par sajtī hai.}
\end{quote}

\begin{quote}
This outfit suits Sara.
\end{quote}

Past (it …): \ur{وہ} \ur{کمرہ} \ur{پھولوں} \ur{سے} \ur{سجا}, \emph{voh kamrā phūloṅ se sajā}.

\begin{cousinweb}
\textbf{\ur{سجنا}} (\emph{sajnā}) is probably related to Hindi \emph{sājnā}, \textquotedblleft{}to be prepared\textquotedblright{}, which CDIAL lists under Sanskrit \emph{sajjyate}, \textquotedblleft{}is equipped\textquotedblright{} (CDIAL 13093); \textbf{\ur{سجانا}}, \textquotedblleft{}to decorate\textquotedblright{}, is its partner.
\end{cousinweb}

\begin{grammarlens}[title={What you've built}]
One more word for getting around.
\end{grammarlens}

\subsection*{Your turn}
\begin{itemize}
\raggedright
  \item \emph{Say it:} \emph{sajnā}
  \item \emph{Say it:} \emph{sajnā}, once more
  \item \emph{Say it:} \emph{yih poshāk sārā par sajtī hai.} — This outfit suits Sara.
  \item \emph{Recall:} say the Urdu for an outfit, then the Urdu for a handkerchief, then say \emph{sajnā} again
\end{itemize}

\begin{tcolorbox}[breakable,colback=teal!4,colframe=teal!35!black,title={Before you move on}]
What does \ur{سجنا} mean? (\textquotedblleft{}To look good (on someone); to be decorated\textquotedblright{}.) Say it once more.
\end{tcolorbox}
